% !TEX root = task-2-1-3.tex
% Editable task bank. Append new BankTask blocks at the end of the document.
\documentclass[12pt,a4paper]{article}
\usepackage{physics-textbook}
\usepackage{graphicx,needspace}
\usetikzlibrary{decorations.pathmorphing}
\newif\ifShowSolutions
\ifdefined\TutorEdition
  \ShowSolutionstrue
\else
  \ShowSolutionsfalse % Change to \ShowSolutionstrue to display solutions.
\fi
\newcounter{banktask}
\newsavebox{\BankTaskBox}
\newcommand{\BankTaskContent}[3]{%
  \noindent\textbf{\thebanktask.}\quad #2\par
  \ifShowSolutions
    \Needspace{50mm}% Keep the solution heading with its opening text or diagram.
    #3\par
    \ifdefined\TutorEdition
      \ifcsname TutorSource#1\endcsname
        \nopagebreak\smallskip{\small\csname TutorSource#1\endcsname\par}%
      \fi
    \fi
  \fi
}
\newcommand{\BankTask}[3]{%
  \refstepcounter{banktask}\par\medskip
  \sbox{\BankTaskBox}{\begin{minipage}{\linewidth}%
    \BankTaskContent{#1}{#2}{#3}\end{minipage}}%
  \ifdim\dimexpr\ht\BankTaskBox+\dp\BankTaskBox\relax>\dimexpr\textheight-20mm\relax
    \Needspace{35mm}\BankTaskContent{#1}{#2}{#3}%
  \else
    \Needspace{\dimexpr\ht\BankTaskBox+\dp\BankTaskBox\relax}%
    \noindent\usebox{\BankTaskBox}\par
  \fi
  \medskip
}
\begin{document}
\BookSubsection{2.1.3}{Взаимодействие частиц вещества}


% Source sdamgia-3205: https://phys-ege.sdamgia.ru/problem?id=3205
\BankTask{sdamgia-3205}{
Какое из утверждений не соответствует представлениям молекулярно-кинетической теории о строении газов?

$1$. Все тела состоят из частиц: атомов, молекул и ионов;

$2$. Частицы находятся в непрерывном хаотическом движении (тепловом);

$3$. Все частицы летают со строго определенными по величине скоростями;

$4$. Частицы взаимодействуют друг с другом путем абсолютно упругих столкновений.
}{
\smallskip\noindent\textbf{Решение}\par
Согласно молекулярно-кинетической теории, все тела, состоят из частиц, которые находятся в непрерывном хаотическом движении и взаимодействуют только посредством упругих столкновений. Естественно, что при столкновениях скорости частиц изменяются. Утверждение $3$ не является верным.

Ответ: $3$.
}

% Source sdamgia-3208: https://phys-ege.sdamgia.ru/problem?id=3208
\BankTask{sdamgia-3208}{
Какие из утверждений справедливы для жидкостей?

А) Характерное расстояние между молекулами жидкости по порядку величины совпадает с размерами самих молекул 

Б) Жидкость сохраняет форму

В) Жидкость сохраняет объем

Г) Молекулы жидкости образуют периодичную решетку

$1$. А и Г 

$2$. Б и Г

$3$. А и В

$4$. Б и В
}{
\smallskip\noindent\textbf{Решение}\par
Молекулы жидкости достаточно плотно упакованы, так что утверждение А верно. В силу плотной упаковки и достаточно сильного взаимодействия между молекулами, жидкости практически несжимаемы. Как следствие, они сохраняют объем. Молекулы жидкости ведут так называемый «кочевой образ жизни», они могут переходить из одного положения в другое, этим обеспечивается текучесть жидкостей, таким образом, жидкости не сохраняют форму. Утверждение Г относится к кристаллическим телам. Окончательно, верными являются утверждения А и В.

Ответ: $3$.
}

% Source sdamgia-3464: https://phys-ege.sdamgia.ru/problem?id=3464
\BankTask{sdamgia-3464}{
В молекулярной физике используется понятие «идеальный газ». Это понятие применимо тогда, когда можно пренебречь:

А. Средней кинетической энергией поступательного движения атомов или молекул.

Б. Энергией взаимодействия атомов или молекул, то есть действием сил притяжения или отталкивания.

В. Массой атомов или молекул.

Какое(-ие) из утверждений правильно(-ы)?

$1$) только А

$2$) только Б

$3$) только В

$4$) А и В
}{
\smallskip\noindent\textbf{Решение}\par
Идеальный газ — это математическая модель газа, в которой предполагается, что потенциальной энергией взаимодействия молекул (или атомов) можно пренебречь по сравнению с их кинетической энергией. В этой модели предполагается, что между молекулами не действуют силы притяжения или отталкивания, соударения частиц между собой и со стенками сосуда абсолютно упруги.

Таким образом, утверждение Б верно, утверждение А неверно. Поскольку кинетической энергией в модели идеального газа не пренебрегают, то не могут пренебрегать и массой частиц газа, так как кинетическая энергия этой массе пропорциональна. Следовательно, для применимости модели идеального газа необходимо выполнение только требования Б.

Ответ: $2$.
}

% Source sdamgia-3543: https://phys-ege.sdamgia.ru/problem?id=3543
\BankTask{sdamgia-3543}{
Какие положения включает модель идеального газа?

А. Объемом, занимаемым молекулами газа, можно пренебречь по сравнению с объемом сосуда.

Б. Взаимодействием молекул можно пренебречь.

В. Молекулы не сталкиваются друг с другом при движении.

$1$. А, Б и В

$2$. А и Б

$3$. А и В

$4$. Б и В
}{
\smallskip\noindent\textbf{Решение}\par
Идеальный газ — математическая модель газа, в которой предполагается, что потенциальной энергией взаимодействия молекул можно пренебречь по сравнению с их кинетической энергией. В этой модели считается, что между молекулами не действуют силы притяжения или отталкивания, то есть взаимодействием молекул пренебрегают. Модель идеального газа хорошо работает только для не очень плотных газов, поскольку в этой модели пренебрегают объемом, занимаемым молекулами газа по сравнению с объемом сосуда. А вот утверждение В к этой модели не относится. Поскольку именно за счет столкновений молекул, в результате которых они обмениваются кинетическими энергиями, газ приходит к состоянию термодинамического равновесия.

Ответ: $2$.
}

% Source sdamgia-4118: https://phys-ege.sdamgia.ru/problem?id=4118
\BankTask{sdamgia-4118}{
В учебнике по физике в одном из абзацев написано: «Молекулы считаются материальными точками, которые хаотически движутся и абсолютно упруго соударяются друг с другом и со стенками сосуда. В промежутках между столкновениями молекулы друг с другом и со стенками сосуда не взаимодействуют». Какая физическая модель описывается в этом абзаце учебника?

$1$) монокристаллическое твердое тело

$2$) поликристаллическое твердое тело

$3$) идеальная жидкость

$4$) идеальный газ
}{
\smallskip\noindent\textbf{Решение}\par
Данное описание относится к модели идеального газа.

Ответ: $4$.
}


% Source sdamgia-6372: https://phys-ege.sdamgia.ru/problem?id=6372
\BankTask{sdamgia-6372}{
Учитель записал на доске три утверждения, относящиеся к молекулам.

$1$. Размерами молекул можно пренебречь.

$2$. Молекулы при столкновениях взаимодействуют как упругие шары.

$3$. При любом расстоянии между молекулами между ними действуют силы притяжения.

Какие из этих утверждений нельзя соотнести с моделью идеального газа?

$1$) все три утверждения

$2$) только третье утверждение

$3$) первое и второе утверждения

$4$) первое и третье утверждения
}{
\smallskip\noindent\textbf{Решение}\par
В модели идеального газа предполагается, что частицы газа малы по сравнению со средним расстоянием между ними, то есть размерами частиц газа пренебрегаем. Частицы взаимодействуют только при столкновения, то есть сил притяжения между ними нет. При этом в столкновениях частицы взаимодействую как упругие шары. Таким образом, верны первое и второе утверждения. То есть утверждение под номером три нельзя соотнести с моделью идеального газа, значит, правильный ответ указан под номером $2$.

Правильный ответ указан под номером $2$.
}

% Source sdamgia-4343: https://phys-ege.sdamgia.ru/problem?id=4343
\BankTask{sdamgia-4343}{
Модель, служащая для демонстрации внутреннего строения тел, устроена следующим образом: в дно прямоугольной коробки воткнуты одинаковые упругие вертикальные стерженьки, на каждый из которых насажен магнитик в виде плоской таблетки. После приведения одного из магнитиков в колебательное движение вскоре начинают хаотически колебаться на стерженьках и все остальные магнитики, отталкиваясь друг от друга. Эта модель лучше всего иллюстрирует поведение молекул

$1$) идеального газа

$2$) жидкости

$3$) твердого тела

$4$) идеального газа и жидкости
}{
\smallskip\noindent\textbf{Решение}\par
Данная модель лучше всего демонстрирует поведение молекул твердого тела, потому что именно в данном агрегатном состоянии частицы совершают малые колебания вокруг положений равновесия, не перемещаясь по объему тела.

Ответ: $3$
}

% Source sdamgia-3286: https://phys-ege.sdamgia.ru/problem?id=3286
\BankTask{sdamgia-3286}{
Броуновское движение частиц пыльцы в воде объясняется

$1$) хаотичностью химических реакций на поверхности частиц

$2$) непрерывностью и хаотичностью теплового движения молекул воды

$3$) существованием сил притяжения и отталкивания между атомами в молекулах

$4$) наличием питательных веществ в воде
}{
\smallskip\noindent\textbf{Решение}\par
Броуновское движение происходит из-за того, что все жидкости и газы состоят из атомов или молекул — мельчайших частиц, которые находятся в постоянном хаотическом тепловом движении, и потому непрерывно толкают броуновскую частицу с разных сторон. Правильный ответ $2$.

Ответ: $2$.
}

% Source sdamgia-5465: https://phys-ege.sdamgia.ru/problem?id=5465
\BankTask{sdamgia-5465}{
Молекулы любого твердого тела находятся в непрерывном движении. Почему тогда тела не рассыпаются на отдельные частицы?

$1$) нет сил отталкивания между молекулами

$2$) среднее значение скорости молекул равно нулю

$3$) действует взаимное притяжение молекул

$4$) действует атмосферное давление
}{
\smallskip\noindent\textbf{Решение}\par
Молекулы твердого тела колеблются в узлах кристаллической решетки вокруг положения равновесия. В этом положении взаимное притяжение молекул сильнее, чем силы отталкивания.

Ответ: $3$.
}

% Source sdamgia-6228: https://phys-ege.sdamgia.ru/problem?id=6228
\BankTask{sdamgia-6228}{
Частицы вещества находятся, в среднем, на таких больших расстояниях друг от друга, при которых силы взаимодействия между ними незначительны. В этом агрегатном состоянии вещество

$1$) хорошо сжимается

$2$) сохраняет и форму, и начальный объем

$3$) сохраняет начальный объем

$4$) сохраняет свою начальную форму
}{
\smallskip\noindent\textbf{Решение}\par
Если частицы вещества находятся на таком расстоянии друг от друга, что взаимодействие между ними можно не учитывать, то в этом агрегатном состоянии вещество хорошо сжимается.

Правильный ответ указан под номером $1$.
}

\end{document}
